\chapter{{\fontsize{14pt}{21pt}\selectfont\MakeUppercase{Анализ предметной области и постановка задачи}}}
\label{cha:analysis}

\section{Анализ предметной области}

Предметная область работы связана с поиском и объяснением сведений из технической документации PostgreSQL. PostgreSQL используется как промышленная система управления базами данных, поэтому документация охватывает широкий набор тем: SQL-команды, типы данных, администрирование, параметры сервера, репликацию, расширения, утилиты и примечания к выпускам~\cite{postgres-docs,date2004,garcia2008,silberschatz2020}. Пользовательские вопросы в такой области редко совпадают с точным названием страницы документации. Обычно вопрос описывает прикладную задачу: найти активные запросы, проверить параметр репликации, объяснить план выполнения, отличить \texttt{VACUUM} от \texttt{ANALYZE}, подобрать источник для дальнейшей ручной проверки.

Техническая документация отличается от справочного текста общего назначения. В ней особенно важны точность, воспроизводимость, связь с конкретным источником и учёт версии. Ошибка в ответе по настройке PostgreSQL может привести не только к неверному пониманию механизма, но и к некорректной конфигурации рабочей базы данных. Поэтому система, формирующая ответ по документации, должна сохранять трассируемость результата: пользователь должен видеть, какие документы и разделы были использованы.

Поиск по документации можно рассматривать как частный случай информационного поиска~\cite{manning2008}. Классические лексические методы хорошо работают, когда запрос содержит точные термины, но хуже обрабатывают переформулировки. Семантический поиск по векторным представлениям лучше переносит разные формулировки, однако сам по себе не гарантирует версионную корректность и не формирует связный ответ. Поэтому в работе рассматривается комбинированное решение: локальный корпус документации, поиск фрагментов, повторное ранжирование, проверка достаточности подтверждений и генерация результата с возвращением источников.

\section{Особенности официальной документации PostgreSQL}

Официальная документация PostgreSQL обладает устойчивой структурой и публикуется по версиям~\cite{postgres-docs,postgres-docs-16,postgres-docs-17,postgres-docs-18}. Для каждой основной версии существует отдельная ветка URL вида \texttt{/docs/<version>/}. Такая организация удобна для пользователя, но требует аккуратной обработки в программной системе: ссылка на страницу текущей версии или на другую ветку документации не должна незаметно подменять выбранную пользователем версию.

Для подготовки корпуса важны несколько свойств документации: заголовки разных уровней, обычные абзацы, списки, таблицы, фрагменты кода и семантические якоря в URL, например название SQL-команды или раздела конфигурации. Эти свойства позволяют сохранять в системе не только текст фрагмента, но и метаданные: заголовок документа, путь раздела, URL источника, тип корпуса и версию.

Официальный корпус в разработанном приложении является основным источником фактов. Дополнительный учебный корпус используется только как вспомогательный слой в обучающем режиме и не должен заменять официальные источники. Такое разграничение соответствует инженерному характеру задачи: объяснение может быть адаптировано для начинающего пользователя, но фактическая часть должна проверяться по документации PostgreSQL.

\section{Проблема учёта версии PostgreSQL}

Политика версионирования PostgreSQL определяет основную версию как первую часть номера версии, начиная с PostgreSQL 10~\cite{postgres-versioning}. Между основными версиями могут изменяться возможности, ограничения, параметры конфигурации, состав утилит и рекомендации по эксплуатации. Поэтому запрос вида <<как настроить логическую репликацию>> не является полностью определённым без указания версии PostgreSQL.

В практической реализации поддерживаются версии PostgreSQL 16, 17 и 18. Они заданы в конфигурации через \texttt{SUPPORTED\_PG\_VERSIONS=16,17,18}, заносятся в таблицу \texttt{versions} и возвращаются маршрутом \texttt{/api/versions}. Версия участвует в валидации входного запроса, выборе документов и дополнительной проверке URL официальных источников. Если пользователь указывает неподдерживаемую версию, запрос отклоняется на уровне схемы данных или при разрешении версии.

Основной риск при отсутствии контроля версии показан в таблице~\ref{tab:analysis-version-risks}. В таблице намеренно рассматриваются не частные ошибки конкретного фрагмента кода, а классы проблем, которые должна предотвращать система.

\begingroup
\centering
\small
\setstretch{1}\selectfont
\setlength{\tabcolsep}{4pt}
\begin{longtable}{|P{0.23\textwidth}|P{0.36\textwidth}|P{0.33\textwidth}|}
\caption{Риски отсутствия учёта версии PostgreSQL}\label{tab:analysis-version-risks}\\
\hline
\tableheadcell{Риск} & \tableheadcell{Проявление} & \tableheadcell{Требуемая реакция системы} \\ \hline
\endfirsthead
\tablecontinuedorend{3}{19}
\hline
\tableheadcell{1} & \tableheadcell{2} & \tableheadcell{3} \\ \hline
\endhead
\endfoot
Смешение веток документации & В контекст одного ответа попадают страницы \texttt{/docs/16/}, \texttt{/docs/17/} и \texttt{/docs/18/}. & Фильтрация по выбранной версии и проверка URL официальных источников. \\ \hline
Подмена текущей версией & Ссылка \texttt{/docs/current/} попадает в ответ для фиксированной версии. & Отсечение официальных источников без ожидаемого фрагмента \texttt{/docs/<version>/}. \\ \hline
Слабый контекст & Найденные фрагменты имеют общие слова, но не подтверждают технический вывод. & Проверка достаточности подтверждений перед генерацией ответа. \\ \hline
Неверное учебное объяснение & Дополнительный материал вытесняет официальный источник в обучающем режиме. & Приоритет официального корпуса; дополнительный корпус только вспомогательный. \\ \hline
\end{longtable}
\par\endgroup

\section{Анализ подходов поиска по технической документации}

Поиск по технической документации может строиться разными способами. Обычный поиск по сайту удобен для навигации, но требует ручной интерпретации страниц. Полнотекстовый поиск использует токенизацию, нормализацию и ранжирование документов; в PostgreSQL такие возможности представлены встроенными механизмами полнотекстового поиска~\cite{postgres-textsearch-16}. Семантический поиск использует векторные представления и расстояние между векторами, что позволяет находить близкие по смыслу фрагменты~\cite{dpr2020,sbert2019}. Комбинированные методы объединяют лексический и семантический сигнал, а повторное ранжирование корректирует порядок кандидатов~\cite{bm25-prf,colbert2020}.

Для задач вопросно-ответного взаимодействия применяется подход генерации с дополненным извлечением: система сначала извлекает контекст из внешнего корпуса, затем передаёт его генеративной модели~\cite{lewis2020rag}. В работе этот подход не трактуется как универсальное средство автоматической истины. Он используется как инженерный способ связать генерацию ответа с локально контролируемыми источниками, ограничить контекст выбранной версией и вернуть пользователю список источников.

Сравнение классов методов приведено в таблице~\ref{tab:analysis-search-approaches}. Из таблицы следует, что ни один из подходов по отдельности не закрывает все требования: связный ответ, источники, локальный контроль корпуса и снижение риска смешения версий достигаются только при сочетании поиска фрагментов, контроля версии и проверки подтверждений.

\begingroup
\centering
\small
\setstretch{1}\selectfont
\setlength{\tabcolsep}{3pt}
\begin{longtable}{|P{0.23\textwidth}|P{0.23\textwidth}|P{0.23\textwidth}|P{0.23\textwidth}|}
\caption{Сравнение подходов к поиску по технической документации}\label{tab:analysis-search-approaches}\\
\hline
\tableheadcell{Подход} & \tableheadcell{Принцип работы} & \tableheadcell{Преимущества} & \tableheadcell{Ограничения} \\ \hline
\tableheadcell{1} & \tableheadcell{2} & \tableheadcell{3} & \tableheadcell{4} \\ \hline
\endfirsthead
\multicolumn{4}{l}{\tablepartlabel{окончание таблицы \thetable}}\\
\hline
\tableheadcell{1} & \tableheadcell{2} & \tableheadcell{3} & \tableheadcell{4} \\ \hline
\endhead
\endfoot
Поиск по ключевым словам & Поиск точных совпадений терминов в тексте документа. & Простота реализации и понятность результата. & Низкая устойчивость к переформулировкам и синонимам. \\ \hline
Семантический поиск & Сравнение векторного представления запроса с векторными представлениями фрагментов. & Лучше обрабатывает близкие по смыслу формулировки. & Может выбрать тематически близкий, но версионно неверный фрагмент без контроля версии. \\ \hline
Полнотекстовый поиск & Индексация текста с нормализацией словоформ и ранжированием документов. & Хорошо подходит для терминологических запросов и локального корпуса. & Не формирует итоговый ответ и не проверяет достаточность источников. \\ \hline
Комбинированный поиск & Объединение семантических и лексических кандидатов. & Повышает устойчивость при технических терминах и свободной формулировке. & Требует логики объединения, фильтрации и повторного ранжирования. \\ \hline
Генерация с извлечением и источниками & Генерация ответа только после извлечения контекста из корпуса. & Формирует связный результат и позволяет вернуть источники. & Качество зависит от корпуса, поиска фрагментов и правил проверки контекста. \\ \hline
\end{longtable}
\par\endgroup

\section{Сравнение существующих инструментов и конкурирующих решений}

Для обоснования направления решения сравнивались не конкретные коммерческие продукты, а классы инструментов, которыми пользователь обычно пользуется при работе с документацией PostgreSQL. Такой подход корректен для инженерной ВКР, поскольку задача состоит не в маркетинговом сравнении, а в выявлении функционального разрыва между доступными способами поиска и требуемым поведением системы.

Критерии сравнения выбраны из требований предметной области: учёт версии, наличие источников, формирование связного ответа, риск смешения версий, локальный контроль корпуса и пригодность для инженерного использования. Результат сравнения приведён в таблице~\ref{tab:analysis-competitors}. В последней строке указано направление, реализованное в практической части проекта.

Обычный поиск по сайту документации PostgreSQL остается полезным при ручной проверке, но он не решает задачу сопровождения ответа источниками в одном пользовательском сценарии. Пользователь получает набор страниц и сам выбирает, какие фрагменты относятся к его вопросу. Для опытного специалиста это приемлемо, однако начинающий пользователь может пропустить важное ограничение версии или принять общий раздел за ответ на частный вопрос. Кроме того, при переходах между страницами легко смешать материалы разных веток документации, особенно если часть ссылок ведет на текущую версию.

Полнотекстовый поиск в локальном корпусе лучше подходит для контролируемой среды, поскольку документы можно заранее разделить по версиям и хранить вместе с метаданными. Его ограничение состоит в том, что терминологически близкий запрос не всегда содержит те же слова, что и документация. Например, пользователь может спрашивать о <<зависших запросах>>, а документация описывает системное представление \texttt{pg\_stat\_activity} и состояние запроса другими терминами. В такой ситуации только лексического совпадения недостаточно: требуется поиск по смысловой близости, а затем дополнительная проверка найденных фрагментов.

Выбор подхода генерации с дополненным извлечением и учётом версии связан именно с этим разрывом между ручным поиском и неконтролируемой генерацией. В разработанном решении генеративная модель не заменяет документацию и не выступает самостоятельным источником знаний. Она получает уже отобранный контекст, ограниченный выбранной основной версией PostgreSQL, а результат возвращается вместе со списком использованных источников. Поэтому в работе важен не сам факт применения генерации, а порядок обработки запроса: сначала валидация версии и поиск фрагментов, затем повторное ранжирование, проверка подтверждений и только после этого формирование ответа.

Такой подход также объясняет разделение официального и дополнительного корпусов. Официальная документация используется как основной источник фактов для краткого и подробного режимов. Дополнительный корпус допустим только в обучающем режиме, где требуется более плавное объяснение и последовательность шагов. При этом дополнительный материал не должен вытеснять официальный источник из ответа: если официального подтверждения недостаточно, система должна отказаться от уверенной формулировки. Это ограничение снижает риск того, что удобное учебное объяснение будет принято за проверенный факт без связи с документацией PostgreSQL.

\begingroup
\centering
\small
\setstretch{1}\selectfont
\setlength{\tabcolsep}{1pt}
\begin{longtable}{|P{0.23\textwidth}|P{0.33\textwidth}|P{0.34\textwidth}|}
\caption{Сравнение инструментов и подходов для поиска по документации PostgreSQL}\label{tab:analysis-competitors}\\
\hline
\tableheadcell{Подход} & \tableheadcell{Возможности} & \tableheadcell{Ограничения для задачи} \\ \hline
\endfirsthead
\tablecontinuedorend{3}{23}
\hline
\tableheadcell{1} & \tableheadcell{2} & \tableheadcell{3} \\ \hline
\endhead
\endfoot
Поиск по сайту документации PostgreSQL & Позволяет вручную выбрать ветку документации и перейти к странице источника. & Не формирует связный ответ, требует ручного анализа и не защищает пользователя от перехода к другой версии. \\ \hline
Полнотекстовый поиск в локальном корпусе & Может работать с локально подготовленными документами и возвращать путь к найденному материалу. & Не формирует итоговый ответ и не проверяет достаточность найденных источников без дополнительной логики. \\ \hline
Общие поисковые системы & Быстро находят страницы, краткие фрагменты выдачи и сторонние материалы по свободной формулировке запроса. & Не гарантируют выбранную версию PostgreSQL, локальный контроль корпуса и инженерную проверяемость ответа. \\ \hline
Локальный HTML-корпус документации & Удобен для опытного пользователя, если локально загружен правильный корпус. & Не формирует ответ и зависит от того, не смешаны ли в локальной папке материалы разных версий. \\ \hline
Вопросно-ответная система без привязки к источникам & Формирует связный текстовый ответ по запросу пользователя. & Может смешивать версии и выдавать неподтверждённые сведения, если источники отсутствуют или неполны. \\ \hline
Генерация с учётом версии и источниками & Использует версию как обязательный параметр, возвращает источники и формирует ответ в выбранном режиме. & Качество результата зависит от полноты корпуса, поиска фрагментов, повторного ранжирования и проверки подтверждений. \\ \hline
\end{longtable}
\par\endgroup

\section{Риски генеративных ответов и необходимость подтверждения источниками}

Генеративная модель может построить грамматически убедительный ответ даже тогда, когда контекст недостаточен. В инженерной задаче это является существенным риском: пользователь может принять неподтверждённую формулировку за факт. Поэтому разработанная система не должна использовать генерацию как самостоятельный источник знаний. Генерация допустима только после извлечения фрагментов из корпуса и проверки, что среди них есть официальные источники выбранной версии.

В приложении эта идея выражается в трёх требованиях. Во-первых, ответ должен сопровождаться списком источников. Во-вторых, в кратком и подробном режимах контекст строится только на официальном корпусе. В-третьих, в обучающем режиме дополнительный корпус может улучшать учебную подачу, но официальный корпус должен оставаться основой фактического вывода. Подобные ограничения соответствуют общим принципам проектирования надёжных информационных систем: ограничения должны быть выражены не только в пользовательской инструкции, но и в программной логике~\cite{martin2017,fowler2002,bass2021}.

Количественная оценка качества ранжирования не входит в состав этой ВКР. Поэтому проверка результата строится функционально и сценарно: подтверждается, что реализованные правила выбора версии, режима, источников и безопасного отказа работают согласно требованиям.

\section{Анализ пользователей и сценариев использования}

Пользователями приложения являются разработчики, администраторы баз данных, инженеры эксплуатации и начинающие специалисты. Разработчику нужны быстрые ответы по SQL-командам, системным представлениям и типичным диагностическим запросам. Администратор баз данных обращается к параметрам сервера, обслуживанию, репликации и ограничениям версии. Инженеру эксплуатации важны воспроизводимые инструкции, контейнерное окружение и история выполненных запросов. Начинающему пользователю нужна пошаговая подача материала и пояснение предварительных условий.

Основные пользовательские сценарии приведены в таблице~\ref{tab:analysis-user-scenarios}. Они служат основой для требований и последующей программы испытаний.

\begingroup
\centering
\small
\setstretch{1}\selectfont
\setlength{\tabcolsep}{4pt}
\begin{longtable}{|P{0.23\textwidth}|P{0.34\textwidth}|P{0.35\textwidth}|}
\caption{Пользовательские сценарии приложения}\label{tab:analysis-user-scenarios}\\
\hline
\tableheadcell{Пользователь} & \tableheadcell{Типовой запрос} & \tableheadcell{Ожидаемое поведение системы} \\ \hline
\endfirsthead
\multicolumn{3}{l}{\tablepartlabel{окончание таблицы \thetable}}\\
\hline
\endhead
\endfoot
Разработчик & <<Напиши SQL-запрос для активных запросов дольше пяти минут>>. & Краткий или подробный ответ с проверяемыми источниками и, если уместно, блоком SQL. \\ \hline
Администратор баз данных & <<Какие параметры проверить для логической репликации в PostgreSQL 16?>> & Поиск официальных источников выбранной версии и отказ от смешения с другими версиями. \\ \hline
Инженер эксплуатации & <<Как проверить блокировки и найти блокирующий запрос?>> & Подробный ответ с источниками и возможностью найти запись в истории. \\ \hline
Начинающий специалист & <<Объясни EXPLAIN ANALYZE простыми словами>>. & Обучающий результат с кратким объяснением, предварительными условиями, шагами и замечаниями. \\ \hline
\end{longtable}
\par\endgroup

\section{Постановка задачи}

Необходимо разработать веб-приложение для работы с документацией PostgreSQL, которое принимает вопрос пользователя, выбранную основную версию PostgreSQL и режим ответа. Система должна подготовить локальный корпус, сохранить документы и фрагменты в базе данных, вычислить векторные представления, выполнить поиск фрагментов, применить контроль версии, выполнить повторное ранжирование, проверить достаточность подтверждений и сформировать ответ с источниками.

Практическая часть ограничивается реализованными функциями. В ней нет пользовательского переключателя корпуса: выбор корпуса определяется режимом ответа. Для краткого и подробного режимов используется официальный корпус. Для обучающего режима используется официальный корпус и дополнительный учебный корпус, причем дополнительный корпус не должен вытеснять официальный.

Задача разработки включает серверную часть, модель данных, маршруты API, предварительную подготовку корпуса, обработку запроса в момент обращения, пользовательский интерфейс и автоматические тесты. Набор поддерживаемых версий фиксируется конфигурацией проекта, а поддержка новых версий предполагает наличие соответствующего корпуса и переиндексацию.

Техническое задание на разработку веб-приложения приведено в приложении~А.

\section{Функциональные требования}

Функциональные требования сформулированы на основе постановки задачи и сверены с фактической реализацией. Требования приведены в приложении~В, таблице~\ref{tab:analysis-fr}. В таблице не перечисляются все внутренние файлы проекта, поскольку требования должны описывать поведение системы, а не структуру репозитория.

\section{Нефункциональные требования}

Нефункциональные требования приведены в таблице~\ref{tab:analysis-nfr}. Они определяют ограничения качества, воспроизводимости и проверяемости, но не вводят неподтверждённых метрик точности.

\begingroup
\centering
\small
\setstretch{1}\selectfont
\setlength{\tabcolsep}{3pt}
\begin{longtable}{|P{0.10\textwidth}|P{0.28\textwidth}|P{0.52\textwidth}|}
\caption{Нефункциональные требования к веб-приложению}\label{tab:analysis-nfr}\\
\hline
\tableheadcell{ID} & \tableheadcell{Требование} & \tableheadcell{Описание} \\ \hline
\tableheadcell{1} & \tableheadcell{2} & \tableheadcell{3} \\ \hline
\endfirsthead
\multicolumn{3}{l}{\tablepartlabel{окончание таблицы \thetable}}\\
\hline
\tableheadcell{1} & \tableheadcell{2} & \tableheadcell{3} \\ \hline
\endhead
\endfoot
NFR-01 & Трассируемость ответа & Ответ должен сопровождаться источниками с названием, URL, типом корпуса, разделом, позицией и оценкой релевантности. \\ \hline
NFR-02 & Версионная изоляция & Официальные источники другой версии или \texttt{/docs/current/} не должны попадать в ответ выбранной версии. \\ \hline
NFR-03 & Воспроизводимость развёртывания & Серверная часть и PostgreSQL должны запускаться через Docker Compose с явной конфигурацией окружения~\cite{docker-docs}. \\ \hline
NFR-04 & Контролируемость генерации & Генерация должна выполняться только по переданному контексту, а источники возвращаются отдельно от текста ответа. \\ \hline
NFR-05 & Тестируемость & Критичные правила должны проверяться автоматическими тестами, запускаемыми командой pytest~\cite{pytest-docs,beizer1990,myers2011}. \\ \hline
NFR-06 & Безопасность конфигурации & Секреты генеративной модели и административный ключ должны передаваться через переменные окружения, а не храниться в коде. \\ \hline
NFR-07 & Расширяемость корпуса & Добавление новой версии должно выполняться через добавление корпуса и переиндексацию без изменения маршрута пользовательского запроса. \\ \hline
NFR-08 & Ограниченность оценки & Оценка результатов должна быть функциональной и сценарной; количественные метрики ранжирования не заявляются без отдельного набора эталонных запросов. \\ \hline
\end{longtable}
\par\endgroup

Нормативная база отчёта включает стандарты СИБИД, ЕСПД и стандарты жизненного цикла автоматизированных систем. Структура пояснительной записки, оформление списка источников, библиографических ссылок, технического задания и критериев качества программного обеспечения соотнесены с этими требованиями~\cite{gost-7-32-2017,gost-r-7-0-100-2018,gost-r-7-0-5-2008,gost-19-201-78,gost-34-601-90,gost-34-602-89,gost-r-iso-25010-2015}.

\section{Выводы}

В первой главе проанализирована область поиска по документации PostgreSQL, выявлена проблема учёта основной версии и рассмотрены подходы к поиску по технической документации. Сравнение показало, что обычный, полнотекстовый и локальный поиск, а также вопросно-ответные системы без источников не обеспечивают одновременно версионную корректность, связность ответа, локальный контроль корпуса и трассируемость источников. На этой основе сформулированы постановка задачи, функциональные и нефункциональные требования к веб-приложению.
